% !TEX root = ../main.tex

\section{Methods}
\label{sec:methods}

Inline math goes between dollar signs, like $E = mc^2$ or $\alpha_i \leq \beta^2$.
Subscripts use \verb|_| and superscripts use \verb|^|; put several characters in braces, as in $x_{\mathrm{max}}$.

A numbered equation uses the \texttt{equation} environment:
\begin{equation}
  y(t) = A \, e^{-t/\tau} \cos(\omega t + \phi)
  \label{eq:model}
\end{equation}
where $A$ is the amplitude, $\tau$ the decay time and $\omega$ the angular frequency.
The \verb|\label| inside lets you write Equation~\ref{eq:model} anywhere in the paper.

Several aligned lines use \texttt{align}; the \verb|&| marks where lines line up and \verb|\\| ends a line:
\begin{align}
  \nabla \cdot \mathbf{E} &= \frac{\rho}{\varepsilon_0} \\
  \nabla \times \mathbf{B} &= \mu_0 \mathbf{J} + \mu_0 \varepsilon_0 \frac{\partial \mathbf{E}}{\partial t}
\end{align}

Comments in OpenTex work like comments in Word: select text, click \textbf{+ Comment}, and a collaborator can reply or resolve it.
They are kept in a separate file in the project folder (\texttt{comments.opentex.json}), so they never clutter the text or appear in the PDF.
